\SourcePage{505}{508}
\section{Interaction between molecular vibration and rotation}

Up to this point, rotation and vibration have been treated as independent motions of a molecule. In fact, when both occur simultaneously, a distinctive interaction arises between them (E.~Teller, L.~Tisza, O.~Placzek, 1932--1933).

\SourcePage{506}{509}
We begin with linear polyatomic molecules. A linear molecule can undergo two kinds of vibration (see the end of §~100): longitudinal vibrations with nondegenerate frequencies, and transverse vibrations with doubly degenerate frequencies. It is the latter that concern us here.

A molecule undergoing transverse vibrations generally has angular momentum. This follows at once from elementary mechanical considerations\footnote{For example, two mutually perpendicular transverse vibrations whose phases differ by $\pi/2$ may be regarded as a pure rotation of the bent molecule about its longitudinal axis.}, but it can also be demonstrated by a quantum-mechanical treatment. Such a treatment also determines the possible values of this momentum in a specified vibrational state.

Suppose that one particular doubly degenerate frequency $\omega_\alpha$ is excited in the molecule. The energy level with vibrational quantum number $v_\alpha$ has degeneracy $v_\alpha+1$. It is represented by $v_\alpha+1$ wave functions
\[
 \psi_{v_{\alpha1}v_{\alpha2}}=\operatorname{const}\cdot
 \exp\left[-\frac12c_\alpha^2(Q_{\alpha1}^2+Q_{\alpha2}^2)\right]
 H_{v_{\alpha1}}(c_\alpha Q_{\alpha1})H_{v_{\alpha2}}(c_\alpha Q_{\alpha2})
\]
(where $v_{\alpha1}+v_{\alpha2}=v_\alpha$), or by any independent linear combinations of them. In every one of these functions, the polynomial multiplying the exponential has the same highest total degree in $Q_{\alpha1}$ and $Q_{\alpha2}$, namely $v_\alpha$. Linear combinations of the functions $\psi_{v_{\alpha1}v_{\alpha2}}$ can evidently always be chosen as basis functions in the form
\begin{equation}
 \begin{split}
 \psi_{v_\alpha l_\alpha}={}&\operatorname{const}\cdot
 \exp\left[-\frac12c_\alpha^2(Q_{\alpha1}^2+Q_{\alpha2}^2)\right]\times\\
 &\times\left[(Q_{\alpha1}+iQ_{\alpha2})^{(v_\alpha+l_\alpha)/2}
 (Q_{\alpha1}-iQ_{\alpha2})^{(v_\alpha-l_\alpha)/2}+\ldots\right].
 \end{split}
 \tag{104.1}
\end{equation}
The square brackets contain a definite polynomial of which only the highest-degree term has been displayed; $l_\alpha$ is an integer having the $v_\alpha+1$ possible values $l_\alpha=v_\alpha$, $v_\alpha-2$, $v_\alpha-4,\ldots,-v_\alpha$.

The normal coordinates $Q_{\alpha1}$ and $Q_{\alpha2}$ of a transverse vibration describe two mutually perpendicular displacements from the molecular axis. Under a rotation through an angle $\varphi$ about this axis, the highest-degree polynomial term, and hence the whole function $\psi_{v_\alpha l_\alpha}$, is multiplied by
\[
 \exp\left\{i\varphi\left(\frac{v_\alpha+l_\alpha}{2}\right)
 -i\varphi\left(\frac{v_\alpha-l_\alpha}{2}\right)\right\}
 =\exp(il_\alpha\varphi).
\]

\SourcePage{507}{510}
It follows that the function (104.1) describes a state whose angular momentum about the axis is $l_\alpha$.

We thus obtain the following result. In a state in which a doubly degenerate frequency $\omega_\alpha$ is excited with quantum number $v_\alpha$, the molecule has an angular momentum about its own axis that takes the values
\begin{equation}
 l_\alpha=v_\alpha,v_\alpha-2,v_\alpha-4,\ldots,-v_\alpha.
 \tag{104.2}
\end{equation}
This is called the molecule's \emph{vibrational angular momentum}. If several transverse vibrations are excited at once, the total vibrational angular momentum is $\sum l_\alpha$. Together with the electronic orbital angular momentum, it gives the total angular momentum $l$ of the molecule about its axis.

As in a diatomic molecule, the total molecular angular momentum $J$ cannot be smaller than the angular momentum about the axis. Thus $J$ assumes the values
\[
 J=|l|,|l|+1,\ldots
\]
In other words, states with $J=0,1,\ldots,|l|-1$ do not exist.

For harmonic vibrations the energy depends only on the numbers $v_\alpha$, not on $l_\alpha$. Anharmonicity removes the degeneracy of the vibrational levels with respect to $l_\alpha$, though only partially: the levels remain doubly degenerate, states related by simultaneous reversal of the signs of every $l_\alpha$ and of $l$ having equal energy. In the approximation following the harmonic one, the energy acquires a term quadratic in the momenta $l_\alpha$, of the form
\[
 \sum_{\alpha,\beta}g_{\alpha\beta}l_\alpha l_\beta
\]
(the $g_{\alpha\beta}$ are constants). The remaining twofold degeneracy is removed by an effect analogous to $\Lambda$-doubling in diatomic molecules.

Before turning to nonlinear molecules, we must make a purely mechanical observation. For an arbitrary nonlinear system of particles, one must ask how vibrational motion can be separated from rotation at all, or, equivalently, what is meant by a ``nonrotating system.'' At first sight, vanishing angular momentum might appear to provide the criterion for the absence of rotation:
\begin{equation}
 \sum m(\mathbf r\times\mathbf v)=0
 \tag{104.3}
\end{equation}

\SourcePage{508}{511}
(the sum is over the particles of the system). The expression on the left, however, is not the total time derivative of any function of the coordinates. Consequently, this equality cannot be integrated with respect to time and recast as the vanishing of some coordinate function. Yet just such a formulation is needed if the concepts of ``pure vibration'' and ``pure rotation'' are to be defined meaningfully.

The condition to be adopted as the definition of no rotation is therefore
\begin{equation}
 \sum m(\mathbf r_0\times\mathbf v)=0,
 \tag{104.4}
\end{equation}
where $\mathbf r_0$ is the radius vector of the equilibrium position of a particle. Writing $\mathbf r=\mathbf r_0+\mathbf u$, with $\mathbf u$ the displacement in a small vibration, gives $\mathbf v=\dot{\mathbf r}=\dot{\mathbf u}$. Integration of (104.4) with respect to time then yields
\begin{equation}
 \sum m(\mathbf r_0\times\mathbf u)=0.
 \tag{104.5}
\end{equation}
We shall regard molecular motion as the combination of a purely vibrational motion satisfying (104.5) and a rotation of the molecule as a whole\footnote{Translational motion is assumed to have been separated from the outset by choosing a coordinate system in which the molecular center of mass is at rest.}. On writing the angular momentum as
\[
 \sum m(\mathbf r\times\mathbf v)=\sum m(\mathbf r_0\times\mathbf v)+\sum m(\mathbf u\times\mathbf v),
\]
we see from the definition (104.4) of no rotation that the sum $\sum m(\mathbf u\times\mathbf v)$ is to be understood as the vibrational angular momentum. It must be remembered, however, that this momentum is only one part of the total momentum of the system and is not itself conserved at all. Thus only a mean value of the vibrational angular momentum can be assigned to each vibrational state.

Molecules having no symmetry axis of order higher than two are of the asymmetric-top type. All vibrational frequencies of molecules of this kind are nondegenerate, since their symmetry groups have only one-dimensional irreducible representations. Hence every vibrational level is nondegenerate. In any nondegenerate state, however, the mean angular momentum vanishes (see §~26). Asymmetric-top molecules therefore have zero mean vibrational angular momentum in every state.

\SourcePage{509}{512}
If the symmetry elements of a molecule include one axis of order higher than two, the molecule is of the symmetric-top type. Such a molecule has both nondegenerate and doubly degenerate vibrational frequencies. The mean vibrational angular momentum again vanishes for the former. A doubly degenerate frequency, on the other hand, corresponds to a nonzero mean projection of angular momentum on the molecular axis.

The energy of rotational motion of a symmetric-top molecule can readily be found with the vibrational angular momentum taken into account. Its operator differs from (103.5) by replacing the top's rotational angular momentum by the difference between the molecule's total conserved angular momentum $\mathbf J$ and its vibrational angular momentum $\mathbf J^{(v)}$:
\begin{equation}
 \widehat H_{\text{rot}}=\frac{\hbar^2}{2I_A}(\widehat{\mathbf J}-\widehat{\mathbf J}^{(v)})^2
 +\frac{\hbar^2}{2}\left(\frac1{I_C}-\frac1{I_A}\right)
 (\widehat J_\zeta-\widehat J_\zeta^{(v)})^2.
 \tag{104.6}
\end{equation}
The energy sought is the mean value $\overline{H}_{\text{rot}}$. Terms in (104.6) containing squares of components of $\mathbf J$ give the purely rotational energy (103.6). Those containing squares of components of $\mathbf J^{(v)}$ give constants independent of the rotational quantum numbers and may be omitted. The terms containing products of components of $\mathbf J$ and $\mathbf J^{(v)}$ represent the effect of the interaction between molecular vibration and rotation that is of interest here. It is called the \emph{Coriolis interaction}, in view of its correspondence to the Coriolis forces of classical mechanics. In averaging these terms, the mean values of the transverse ($\xi,\eta$) components of the vibrational angular momentum must be taken to vanish. The resulting mean Coriolis-interaction energy is
\begin{equation}
 E_{\text{Cor}}=-\frac{\hbar^2}{I_C}kk_v,
 \tag{104.7}
\end{equation}
where the integer $k$, as in §~103, is the projection of the total angular momentum on the molecular axis, while $k_v=\overline{J_\zeta^{(v)}}$ is the mean projection of the vibrational angular momentum characterizing the given vibrational state. Unlike $k$, $k_v$ need not be an integer at all.

Finally, consider spherical-top molecules. These are molecules possessing the symmetry of one of the cubic groups. They have nondegenerate, doubly degenerate, and triply degenerate frequencies, corresponding to the occurrence of one-, two-, and three-dimensional irreducible representations among those of the cubic groups. As always, anharmonicity partially removes the degeneracy of the vibrational levels;

\SourcePage{510}{513}
after these effects have been allowed for, only nondegenerate, doubly degenerate, and triply degenerate levels remain. We shall discuss precisely these levels after their splitting by anharmonicity.

For a spherical-top molecule, the mean vibrational angular momentum plainly vanishes not only in nondegenerate vibrational states but also in doubly degenerate ones. This follows directly from elementary symmetry considerations. Indeed, the vectors of the mean momenta in the two states belonging to one degenerate energy level would have to transform into each other under every symmetry transformation of the molecule. No cubic symmetry group, however, permits two directions that transform only into one another; only sets containing at least three directions are transformed among themselves.

The same argument shows that the mean vibrational angular momentum is nonzero in states belonging to triply degenerate vibrational levels. After averaging over the vibrational state, this momentum is represented by an operator whose matrix elements correspond to transitions among the three mutually degenerate states. In accordance with their number, this operator must have the form $\zeta\widehat{\mathbf I}$, where $\widehat{\mathbf I}$ is the operator of an angular momentum equal to unity (for which $2I+1=3$), and $\zeta$ is a constant characteristic of the given vibrational level. After this averaging, the Hamiltonian of the rotational motion,
\[
 \widehat H_{\text{rot}}=\frac{\hbar^2}{2I}(\widehat{\mathbf J}-\widehat{\mathbf J}^{(v)})^2,
\]
becomes the operator
\begin{equation}
 \widehat H_{\text{rot}}=\frac{\hbar^2}{2I}\widehat{\mathbf J}^{\,2}
 +\frac{\hbar^2}{2I}\overline{\widehat{\mathbf J}^{(v)2}}
 -\frac{\hbar^2}{I}\zeta\widehat{\mathbf J}\widehat{\mathbf I}.
 \tag{104.8}
\end{equation}
The eigenvalues of the first term give the usual rotational energy (103.4), while the second term contributes an inessential constant independent of the rotational quantum number. The last term in (104.8) gives the required Coriolis splitting energy of the vibrational level. The eigenvalues of $\mathbf J\mathbf I$ are obtained in the usual way. For a specified $J$, it can have three distinct values, corresponding to values $J+1$, $J-1$, and $J$ of the vector $\mathbf I+\mathbf J$:
\begin{equation}
 E_{\text{Cor}}^{(J+1)}=-\frac{\hbar^2}{I}\zeta J,\qquad
 E_{\text{Cor}}^{(J-1)}=\frac{\hbar^2}{I}\zeta(J+1),\qquad
 E_{\text{Cor}}^{(J)}=\frac{\hbar^2}{I}\zeta.
 \tag{104.9}
\end{equation}
